% Cómo usarlo
% Se debe colocar los archivos de la portada dentro de una carpeta portada
% desde donde será llamado.
%
% Se deben declarar los siguientes comandos, ajustando los valores por
% los que guste:
%
% \newcommand{\HBar}{\rule{\linewidth}{0.5mm}}
% \newcommand{\logoURL}{logoURL}
% \newcommand{\nAutor}{Nombre del autor}
% \newcommand{\nCurso}{Curso}
% \newcommand{\cTituloA}{Titulo 1}
% \newcommand{\cTituloB}{Titulo 2}
% \newcommand{\cFecha}{Agosto de 2,014.}

%\pagenumbering{gobble}
%\begin{titlepage}

%\vspace*{\fill}
\begin{center}
% logo URL
\includegraphics[width=0.5\textwidth]{latex-url-portada/\logoURL}~
\\[8mm]
\text{\LARGE Universidad Rafael Landívar}
\\
\text{\Large Facultad de Ingeniería}
\\
\text{\large Ingeniería en Informática y Sistemas}
\\[1cm]
\text{\large \nCurso} \\
\nAutor
\\[0.5cm]
\HBar \\[0.5cm]
\text{\LARGE\cTituloA}
\\[0.5cm]
\text{\LARGE\cTituloB} \\[0.5cm]
\HBar \\[1cm]
\cFecha\\[3mm]
\begin{minipage}{0.5\textwidth}
\centering 
\includegraphics{latex-url-portada/img/by-sa}\\
\tiny Este obra está bajo una \href{http://creativecommons.org/licenses/by-sa/4.0/}{licencia de Creative Commons Reconocimiento-CompartirIgual 4.0 Internacional}.
\end{minipage}
\end{center}
%\vspace*{\fill}
%\end{titlepage}
%\pagenumbering{arabic}
